% v20: current frontier and reusable next steps; stable support identifiers.
\subsection{残る五型と，再開する具体的な入口}
\label{ns20:support-agenda}

現在の高分母 support の必要候補は\emph{五型・六十六標識}である．
$G_{5,11}$ は分母上限を用いず全排除され，その六十標識は候補から外れる．
再開時は最初に \nolinkurl{00_CONTINUATION_POLICY_ja.txt} を読む．
以下の計画は，面白さとインパクトを最大化するために柔軟に変更する．

全候補は五頂点を使い，最小次数二以上，二点連結，辺数七以上，$b_1\ge3$ である．
各二価頂点の二つの高近傍はともに次数四であり，その二価頂点を含む高分母三角形を作る．
二価頂点数は零・一・三に限られ，一個の場合は $G_{5,15}$，
三個の場合および七高辺の場合は book $G_{5,13}$ である．
一つの core に高分母 ray が高々一本なら全分母一となる．
表の $ij$ は無向辺，$N$ は標識付き個数であり，補 graph では孤立頂点を省略する．
\begin{center}
\footnotesize
\setlength{\tabcolsep}{4pt}
\begin{tabular}{crclcr}
型&$|E|$&次数列&$E$&低辺の補 graph&$N$\\\hline
$G_{5,13}$&7&$(44222)$&$\{01,02,03,04,12,13,14\}$&$K_3$&10\\
$G_{5,14}$&8&$(43333)$&$K_5-\{14,23\}$&$2K_2$&15\\
$G_{5,15}$&8&$(44332)$&$K_5-\{24,34\}$&$P_3$&30\\
$G_{5,16}$&9&$(44433)$&$K_5-\{34\}$&$K_2$&10\\
$G_{5,17}$&10&$(44444)$&$K_5$&$\varnothing$&1
\end{tabular}
\end{center}
全 $1024$ 標識付き graph の検算は
\nolinkurl{verify_support_agenda_v20.py} と
\nolinkurl{support_agenda_verification_v20.json} にある．
候補の個数は整数接着や幾何的 atlas の存在数ではない．

\paragraph{第一候補：book の混合枝と二つの向き．}
$G_{5,13}$ の二価 cores $2,3,4$ は全て高近傍 $0,1$ を持つ．
低三角形の全正枝・全負枝は排除されたので，負枝の個数は一または二に限られる．
全ての低 $D$ は非零であり，各 core で一正一負となって有向三角形を作る．
一方，高 ray $0$ への正 entry が定める向きも有向三角形となる．
この二つの向きは定義が異なり，一致を仮定しない．
まず両向きが一致する場合と逆の場合に分け，
$L\in\{1,3,4\}$ と有限な node offset を同時に課す．
固定 core の完全有限完成法を用い，混合枝で左右の相手が同時に存在する条件を比較する．
全負枝に対して証明された $w_1\le16$ を，混合枝に流用しない．
低三角形自体が残った場合に中心 $0,1$ の row/pair，$q_{01}$，
高辺の node と $D$ を加える．

\paragraph{並行候補：一つの二価 core と密な支持．}
$G_{5,15}$ は二価 core を一つだけ持つ．
その高近傍はともに四価なので，$G_{5,11}$ の三価相手による上界をそのまま移さない．
共通低 puncture $(23;4)$ の条件と，三高 matching の固定低 row/column 完全復元を組み合わせる．
$G_{5,14}=K_{2,2,1}$ と $G_{5,16},G_{5,17}$ では，
三高分母と四高分母の条件を分け，四高側へ matching 定理を無条件に拡張しない．
高分母 support $H$ と素数 support $G_\ell$ を区別する．
dense mod 五の二模型の整数 lifting も残し，少なくとも二辺の五進付値一を保つ．

\paragraph{局所符号問題の残部．}
三高 $D$ の同時非負性を調べる範囲は
$b=-\beta\ge2$，$d=-D_0\ge1$，$\alpha\ge k\ge b+1$ であり，
$\Delta\ge(-\mathcal B)(-\mathcal B-1)>0$ となる．
特異 Schur 系はこの符号問題では再探索しない．
残る非特異領域では Cramer 分子の整除性，$S,T,U$ の正整数性，
canonical entries の正値性を同時に使う．
特異な正整数局所模型は一般分類の資料として保持し，その存在と同時非負性を区別する．

\paragraph{再利用と，解決済みの型．}
$K_4$，二葉三角形，$G_{5,5},G_{5,7},G_{5,9},G_{5,10},G_{5,11},G_{5,12}$ は再探索しない．
$G_{5,11}$ の証明で用いた，reverse data の解析的上界，
scale の点・全半直線表示，分母一致と $D$ の奇偶障害は他の型にも再利用できる．
ただし上界の各向きで，相手の高次数と例外接着の仮定を確認する．
$G_{5,15},G_{5,16},G_{5,17}$ は $K_4$ を部分 graph に持つが，
support 全体が $K_4$ である場合の排除はそこへ移せない．
整数候補が残れば actual weights，basin areas，integer framings，作用切片へ進む．
幾何的同値・別退化・局所置換へ研究の主題を移すことも，成果の価値に応じて判断する．
